\section{Adaptive precision and two maximal depths at every large prime}
\label{adp-section}

Let \(n>324\) be prime, \(B=n^4\), and \(L=\log(8B)\). For integers
\(B\leq k<2B\), put
\[
 w_k=3k+\zeta,\quad Q_k=N(w_k),\quad
 \alpha_k=w_k/\bar w_k,\qquad \zeta^2=\zeta-1.
\]
Fix a subset \(I\) of this actual block on which the ratios
\(\alpha_k\) are multiplicatively independent. The private-norm
class of Theorem~\ref{pn-private}, or one fixed color of
Theorem~\ref{ncs-colors}, supplies such a set. For
\(w_k^n=A_k+C_k\zeta\), write
\[
 T_k=|A_kC_k(A_k+C_k)|,\qquad
 c_k=\max(|A_k|,|C_k|,|A_k+C_k|),\qquad t_k=\log c_k .
\]
The actual powers are primitive and unramified. Their three nonzero
boundary arms are therefore pairwise coprime.

\begin{lemma}[The individual-arm cap]\label{adp-arm-cap}
For every \(k\in I\),
\[
 n\log(3B)\leq t_k<nL.
\]
For every prime \(q>8B\) dividing \(T_k\),
\begin{equation}\label{adp-depth-cap}
 v_q(T_k)\log q\leq t_k<nL.
\end{equation}
Its homogeneous rank is \(n\), and \(q\equiv1\) or \(-1\pmod{3n}\).
\end{lemma}
\begin{proof}
The norm satisfies \(9B^2\leq Q_k<36B^2\). The quadratic-coordinate
bound and the largest-arm norm inequality give
\[
 Q_k^{n/2}\leq c_k\leq\frac2{\sqrt3}Q_k^{n/2}
       <\frac2{\sqrt3}(6B)^n<(8B)^n.
\]
The last strict inequality follows from \(2/\sqrt3<4/3\leq(4/3)^n\).
At a prime dividing \(T_k\), exactly one arm is divisible by that
prime; its entire multiplicity is bounded by the logarithm of that
arm. This proves \eqref{adp-depth-cap}. Such a prime does not divide
the norm, since the norm reduces to a nonzero coordinate square on
each arm. The established homogeneous rank formula gives rank
dividing \(n\). Rank one would divide the old boundary
\(3k(3k+1)\), which is impossible for \(q>8B>3k+1\). Thus the rank
is \(n\). The split/inert norm-one order formula gives
\(3n\mid q-1\) or \(3n\mid q+1\), respectively.
\end{proof}

\begin{lemma}[Separate moment length and precision]\label{adp-precision}
Let \(q>8B\) be prime, and let \(e,\nu\) be positive integers with
\[
 e\log q\geq2\nu L.
\]
If \(M_e=\#\{k\in I:q^e\mid T_k\}>0\), then
\begin{equation}\label{adp-choose}
 \binom{M_e+\nu-1}{\nu}\leq3n.
\end{equation}
An empty hit set has zero positive-length multisets.
\end{lemma}
\begin{proof}
The proof of Lemma~\ref{um-rigidity} bounds the integer determinant
of any two actual \(\nu\)-fold root products by
\[
 |D|\leq8\nu\,7^{2\nu-1}B^{2\nu-1}<(8B)^{2\nu}.
\]
For completeness, division by \((8B)^{2\nu}\), with \(x=49/64\),
bounds it by
\[
 \frac8{7B}\nu x^\nu
 <\frac8{7B}\sum_{j=0}^{\nu-1}x^j
 <\frac{512}{105B}<1.
\]
The hypothesis gives \(q^e\geq(8B)^{2\nu}\). Equal ratio products
modulo \(q^e\) imply \(q^e\mid D\), hence \(D=0\) and exact ratio
equality in \(\mathbb Q(\zeta)\).

At every hit the conjugate denominator is a unit. The cubic boundary
identity implies that the ratio is killed by \(3n\) modulo \(q^e\).
The split/inert simple-lifting result of Lemma~\ref{mc-torsion}
bounds that norm-one torsion group by \(3n\), since \(q\nmid3n\).
The multiset product map lands in this group. Equality of images
gives exact equality of algebraic ratio products; independence then
recovers all multiplicities. The map is injective, and stars and
bars proves \eqref{adp-choose}. Repeated indices are included.
\end{proof}

\begin{theorem}[A complete bound at each prime]\label{adp-prime}
For every prime \(q>8B\),
\begin{align}
 \sum_{k\in I}(v_q(T_k)-3)_+\log q&\leq13nL,
                    \label{adp-thirteen}\\
 \frac1B\sum_{k\in I}\frac{(v_q(T_k)-3)_+\log q}{t_k}
       &\leq\frac{26}{B}. \label{adp-twentysix}
\end{align}
There is no upper restriction on \(q\).
\end{theorem}
\begin{proof}
Put \(r=\lceil\sqrt n\rceil\) and
\[
 h=\max\left(4,\left\lceil\frac{2rL}{\log q}\right\rceil\right).
\]
Lemma~\ref{adp-precision} with \((e,\nu)=(4,2)\) gives
\(M_4\leq\sqrt{6n}\). With \((e,\nu)=(h,r)\), it gives \(M_h\leq3\).
Indeed, \(M_h\geq4\) would imply
\[
 3n\geq\binom{r+3}{3}>\frac{r^3}{6}
       \geq\frac{n^{3/2}}6>3n,
\]
using \(n>324\).

If \(h=4\), at most three indices have any positive excess, and
\eqref{adp-depth-cap} pays their entire depths by \(3nL\).
Otherwise \(h\) is the indicated ceiling and
\((h-4)\log q\leq2rL\). The full layers \(4,\ldots,h-1\) cost at most
\[
 (h-4)M_4\log q\leq2r\sqrt{6n}L
       \leq4\sqrt6\,nL<10nL.
\]
Every layer from \(h\) onward is supported on at most three indices;
paying their full depths adds at most \(3nL\). The identity
\((e-3)_+=\sum_{j\geq4}\mathbf1_{j\leq e}\) accounts for all layers.
Subtracting four absorbs the ceiling error, without an extra
\(\log q\) term. This proves \eqref{adp-thirteen}.
Now use \(t_k\geq n\log(3B)\) and \(L\leq2\log(3B)\).
\end{proof}

If the actual support is partitioned into \(D\) independent colors
chosen before \(q\), both bounds are multiplied by \(D\).
For a real \(Z>8B\), the same Brun--Titchmarsh input as
Theorem~\ref{um-mean} counts the two eligible progressions and gives
\begin{equation}\label{adp-window}
 \frac1B\sum_{k\in I}
  \frac{\sum_{8B<q\leq Z}(v_q(T_k)-3)_+\log q}{t_k}
 \leq\frac{52Z}{B(n-1)\log(Z/(3n))}
 \leq\frac{55Z}{Bn\log(Z/(3n))}.
\end{equation}
Here \(\varphi(3n)=2(n-1)\), and \(52n/(n-1)<55\) for \(n>324\).
This remains true for \(Z>B^2\), although its right side is generally
too large there. On the established logarithmic window it improves
a constant, rather than the asymptotic range.

For a real \(Z\geq8B\), define the actual finite set
\[
 \mathcal D_Z(I)=\{q>Z:q\text{ prime and }v_q(T_k)\geq4
                         \text{ for some }k\in I\}.
\]
Every label divides the nonzero integer \(\prod_{k\in I}T_k\).
Summing \eqref{adp-twentysix} over just these labels proves
\begin{equation}\label{adp-far-card}
 \frac1B\sum_{k\in I}
  \frac{\sum_{q>Z}(v_q(T_k)-3)_+\log q}{t_k}
 \leq\frac{26|\mathcal D_Z(I)|}{B}.
\end{equation}
The sufficient condition \(|\mathcal D_Z(I)|=o(B)\) has not been
proved. Counting roots at a fixed prime does not count different
primes owned by different roots. This positive target is also
stronger than a signed target that can use depth-one and depth-two
credits; no equivalence is claimed.

\begin{theorem}[Two maximal depths and the complete remainder]
\label{adp-two-owners}
For a prime \(q>8B\), put \(e_k=v_q(T_k)\) and
\[
 r=\lceil\sqrt{6n}\rceil,\qquad
 h_0=\max\left(4,\left\lceil\frac{6L}{\log q}\right\rceil\right),
 \qquad
 h_1=\max\left(h_0,\left\lceil\frac{2rL}{\log q}\right\rceil\right).
\]
Let \(O_q\) be the \(\min(2,|I|)\) indices of largest actual depth,
with ties broken by index order. Then
\begin{align}
 \sum_{k\in I\setminus O_q}(e_k-3)_+\log q
       &\leq39n^{5/6}L, \label{adp-remainder}\\
 \sum_{k\in I}(e_k-3)_+\log q
       &\leq(2n+39n^{5/6})L. \label{adp-two-full}
\end{align}
The normalized remainder is at most \(78n^{-1/6}/B\).
\end{theorem}
\begin{proof}
At precisions \(4,h_0,h_1\), respectively, Lemma~\ref{adp-precision}
gives
\[
 M_4\leq\sqrt{6n},\qquad M_{h_0}<3n^{1/3},\qquad M_{h_1}\leq2.
\]
The middle bound follows from
\(M^3\leq6\binom{M+2}{3}\leq18n<27n\).
For the last bound, \(M\geq3\) would give
\[
 3n\geq\binom{r+2}{2}
       =\frac{(r+1)(r+2)}2>\frac{r^2}2\geq3n.
\]
Every \(h_1\)-deep index belongs to \(O_q\): an omitted such index
and the two selected indices would contradict \(M_{h_1}\leq2\).
Empty and singleton sets satisfy this directly.

For each remaining depth \(e<h_1\), the exact finite inequality is
\begin{equation}\label{adp-pointwise}
 (e-3)_+\leq(h_0-4)\mathbf1_{e\geq4}
                   +(h_1-h_0)\mathbf1_{e\geq h_0}.
\end{equation}
It follows by splitting \(e<4\), \(4\leq e<h_0\), and
\(h_0\leq e<h_1\). Both ceilings satisfy
\[
 (h_0-4)\log q\leq6L,\qquad
 (h_1-h_0)\log q\leq2rL,
\]
including an empty range. Summing \eqref{adp-pointwise} bounds the
complete remainder by
\[
 (6M_4+2rM_{h_0})L
 \leq15\sqrt n\,L+24n^{5/6}L
 \leq39n^{5/6}L,
\]
where \(r\leq\sqrt{6n}+1\leq4\sqrt n\).
The at most two removed indices cost at most \(2nL\) by the
individual-arm cap. Normalizing as before proves the last assertion.
\end{proof}

If \(\log q\geq(r/2)L\), then \(h_0=h_1=4\), so every such prime
with positive excess has at most two hit roots. Its actual
prime-labelled incidence graph allows singleton edges and repeated
edges with different prime labels. The choices \(O_q\) vary with
\(q\). Neither this graph nor the uniform remainder per prime
bounds their unrestricted cumulative costs.

\paragraph{Finite formal scope.}
The separate module
\texttt{Adaptive\allowbreak Owner\allowbreak Arithmetic.lean}
has twelve freshly compiled declarations. They prove the cubic and
two-index binomial thresholds, actual natural ceiling bounds,
selection of two maxima from a finite set, deep-set containment,
\eqref{adp-pointwise} and its complete finite sum, the exact weighted
remainder \((6M_4+2rM_{h_0})L\), and addition of at most two full caps.
The final theorem derives a genuine selected set \(O\) and both
finite budgets from explicit real-height and symmetric-count
interfaces. Natural truncated subtraction is performed before the
cast to real numbers. The old fifteen local declarations were also
freshly compiled, with the pinned Mathlib cache reused; all axiom
queries use only the standard three axioms.

The arithmetic construction of valuations, private witnesses,
torsion and determinant rigidity remains ordinary mathematics here.
The constants \(39,78\) involving fractional real powers and the
Brun--Titchmarsh consequences are not claimed as new Lean theorems.
The roots outside \(I\), pointwise exceptions, the changing maximal
depths across different primes, and the full signed far tail remain
unresolved.
